\section*{Capítulo \ref{ch:b3:endocrinology} --- Endocrinologia e homeostase}

\begin{solution}{exo:b3:endocrinology:1}
\omterm{def:b3:endocrinology:glucose}{Insulina}: peptídeo, receptor de superfície (uma tirosina quinase), age
em minutos, meia-vida de cerca de \qty{5}{min}. Cortisol: esteroide,
\omterm{def:b3:endocrinology:hormone}{receptor nuclear}, horas, meia-vida de \qty{80}{min} (em boa parte ligado
a um transportador). Adrenalina: amina, receptores de superfície
acoplados a proteína~G, segundos, meia-vida de cerca de \qty{2}{min}.
Tiroxina: aminoácido iodado, \omterm{def:b3:endocrinology:hormone}{receptor nuclear}, dias, meia-vida de
\qty{7}{d} (ligada a transportadores). Vasopressina: peptídeo,
receptor de superfície acoplado a proteína~G (cAMP no \omterm{def:b3:renal-osmoregulation:nephron}{ducto coletor}),
minutos, meia-vida de cerca de \qty{15}{min}.
\end{solution}

\begin{solution}{exo:b3:endocrinology:2}
\omterm{def:b3:endocrinology:axes}{Hipotálamo} (TRH) $\to$ \omterm{def:b3:endocrinology:axes}{hipófise} (TSH) $\to$ tireoide (T$_{4}$), com a
T$_{4}$ inibindo os dois níveis acima. Tireoide destruída: T$_{4}$
baixa, TSH alto. \omterm{def:b3:cancer-biology:hallmarks}{Tumor} secretor de TSH: TSH alto e T$_{4}$ alta (a
retroalimentação não consegue suprimir o \omterm{def:b3:cancer-biology:hallmarks}{tumor}). Excesso de comprimidos
de tiroxina: T$_{4}$ alta, TSH suprimido. Carência de iodo: T$_{4}$
baixa ou no limite inferior, TSH alto, e a glândula cresce sob ele num
bócio.
\end{solution}

\begin{solution}{exo:b3:endocrinology:3}
O encéfalo morre de glicose baixa em minutos e só é lesado por uma
glicose alta ao longo de anos; e a evolução encontrou a fome muito mais
vezes que os banquetes. Assim, a defesa contra uma queda é redundante
(\omterm{def:b3:endocrinology:glucose}{glucagon}, adrenalina, cortisol, \omterm{def:b3:endocrinology:hormone}{hormônio} do crescimento) e a defesa
contra uma elevação é única. Previsão: \omterm{def:b3:endocrinology:glucose}{insulina} demais é letal de modo
agudo (coma hipoglicêmico), e de menos é uma doença crônica --- que é o
que a clínica vê.
\end{solution}

\begin{solution}{exo:b3:endocrinology:4}
Um \omterm{def:b3:endocrinology:clock}{zeitgeber} é uma pista ambiental que sincroniza um relógio: a luz, o
horário das refeições, a temperatura, o exercício, a agenda social. O
relógio-mestre usa a luz, pelas \omterm{def:b3:sensory-systems:retina}{células ganglionares} de melanopsina da
\omterm{def:b3:sensory-systems:retina}{retina}. Depois de oito fusos horários o relógio se desloca apenas cerca
de uma hora por dia, de modo que, por uma semana, o sono, o cortisol, a
temperatura e o apetite rodam no horário antigo enquanto os relógios
periféricos, reajustados pelas refeições, derivam em seu próprio ritmo:
o jet lag, pior rumo ao leste, já que adiantar o relógio é mais difícil
que atrasá-lo.
\end{solution}

\begin{solution}{exo:b3:endocrinology:5}
Resposta de metade do máximo com $250$ receptores ocupados: $H = K_{d}\times
250/(\num{20000} - 250) = \qty{1}{nM}\times 0.0127 = \qty{12.7}{pM}$, oitenta vezes abaixo de $K_{d}$. Com \num{2000}
receptores: $250/1750
\times \qty{1}{nM} = \qty{143}{pM}$, onze vezes
menos sensível. Uma \omterm{def:b3:endocrinology:glucose}{insulina} alta e prolongada regula negativamente seus
receptores, de modo que um dado efeito exige mais \omterm{def:b3:endocrinology:glucose}{insulina}: um
componente da resistência à \omterm{def:b3:endocrinology:glucose}{insulina}, que se alimenta de si mesma.
\end{solution}

\begin{solution}{exo:b3:endocrinology:6}
$T_{4} = 13$: $1.5\,\mathrm{e}^{0.92} = \qty{3.8}{mU/L}$; $11$:
$1.5\,\mathrm{e}^{1.84} = \qty{9.4}{mU/L}$; $9$: $1.5\,\mathrm{e}^{2.76}
= \qty{24}{mU/L}$; $25$: $1.5\,\mathrm{e}^{-4.6} = \qty{0.015}{mU/L}$. Em
\qty{11}{pmol/L} a T$_{4}$ ainda está dentro de \qtyrange{10}{20}{}
enquanto o TSH é mais que o dobro de seu limite superior (hipotireoidismo
subclínico); em $13$ o TSH está no limite da faixa; em $9$ os dois estão
anormais; em $25$ os dois estão anormais no outro sentido.
\end{solution}

\begin{solution}{exo:b3:endocrinology:7}
$L = s\beta/(a\gamma) = 0.2\times 0.05/(0.02\times 0.1) = 5$. Excesso
estacionário $g^{*} = (u/a)/(1 + L) = (2/0.02)/6 = \qty{16.7}{mg/dL}$;
sem retroalimentação, $u/a = \qty{100}{mg/dL}$. Excesso de \omterm{def:b3:endocrinology:glucose}{insulina}
$i^{*} = \beta g^{*}/
\gamma = 0.05\times 16.7/0.1 = \qty{8.3}{mU/L}$.
\end{solution}

\begin{solution}{exo:b3:endocrinology:8}
Discriminante $(a - \gamma)^{2} - 4s\beta = 0.0064 - 0.04 < 0$:
oscilatório. $\omega = \sqrt{a\gamma + s\beta - (a+\gamma)^{2}/4} =
\sqrt{0.002 + 0.01 - 0.0036} = \qty{0.092}{min^{-1}}$, período $2\pi/
\omega = \qty{68}{min}$; a amplitude cai por $\mathrm{e}$ em $2/(a +
\gamma) = \qty{16.7}{min}$. Reduzindo $s$ à metade: $s\beta = 0.005$,
discriminante $0.0064 - 0.02 < 0$, ainda oscilatório, $\omega = \sqrt{0.007 - 0.0036}
= \qty{0.058}{min^{-1}}$, período \qty{108}{min};
o tempo de decaimento não muda; $L = 2.5$. Uma alça mais fraca retorna
mais devagar e sustenta um erro estacionário maior.
\end{solution}

\begin{solution}{exo:b3:endocrinology:9}
Por um ano a prednisolona suprimiu o CRH e o ACTH; o córtex da adrenal,
sem estímulo, atrofiou. Ao parar não há esteroide exógeno, o \omterm{def:b3:endocrinology:axes}{hipotálamo}
e a \omterm{def:b3:endocrinology:axes}{hipófise} levam semanas para retomar, e a adrenal encolhida não
conseguiria responder ao ACTH se ele viesse. Uma infecção exige dez
vezes o cortisol normal; sem nenhum, os vasos dilatam e a pressão e a
glicose caem: crise adrenal. Medidos: ACTH baixo (a lesão é central),
cortisol baixo e uma resposta pobre do cortisol ao ACTH injetado (a
glândula se atrofiou) --- ao contrário da doença de Addison, em que o
ACTH é alto. Daí a retirada lenta.
\end{solution}

\begin{solution}{exo:b3:endocrinology:10}
Uma variável: $\dot x = f(x)$ com $f' < 0$ tem um único ponto fixo
$x^{*}$; para $x > x^{*}$, $\dot x < 0$ e $x$ decresce monotonicamente
rumo a $x^{*}$, sem nunca cruzá-lo (unicidade das soluções), e
simetricamente abaixo: nenhuma oscilação. Duas variáveis: a matriz pode
ter autovalores complexos, como no teorema, quando o efetor age por uma
segunda variável com sua própria escala de tempo. Uma repressão
instantânea se acomodaria num nível estável; o atraso entre a
transcrição e a repressão nuclear (tradução, dimerização, fosforilação,
importação) deixa a proteína ultrapassar o alvo antes de o gene ser
desligado e ficar aquém antes de ele ser religado, e o período é
aproximadamente o dobro da soma do atraso com o tempo de degradar a
proteína. Uma degradação mais rápida da PER encurta o período: uma
mutação desestabilizadora da PER2 humana dá um relógio de cerca de 20
horas e uma família que adormece às 7 da noite.
\end{solution}

\begin{solution}{exo:b3:endocrinology:11}
Expoente $k(\text{Ca} - 1.2)$: em $1.10$, $\mathrm{e}^{-2}$, PTH $= 1/
(1 + 0.135) = 0.88$; $1.15$: $0.73$; $1.20$: $0.50$; $1.25$: $1/(1 +
2.72) = 0.27$; $1.30$: $1/(1 + 7.39) = 0.12$. Inclinação no \omterm{thm:b3:endocrinology:loop}{ponto de
ajuste} $-k/4 = -5$ por mmol/L: \qty{5}{\%} do máximo a cada
\qty{0.01}{mmol/L}. O cálcio precisa ser mantido dentro de alguns por
cento e seu efetor age sobre um tampão enorme (o osso) sem ultrapassar o
alvo, de modo que uma resposta íngreme é segura e necessária. Uma alça
de glicose assim íngreme, agindo por uma \omterm{def:b3:endocrinology:glucose}{insulina} com seu próprio tempo
de depuração, teria um \omterm{thm:b3:endocrinology:loop}{ganho de alça} grande e oscilaria fundo na
hipoglicemia depois de cada refeição.
\end{solution}

\begin{solution}{exo:b3:endocrinology:12}
A glicação é irreversível e lenta, de modo que uma célula de idade
$\tau$ tem uma fração $kG\tau$ glicada; em média sobre células com idades
uniformes de \numrange{0}{120} dias, a HbA1c $= kG\times\qty{60}{d}$,
proporcional à glicose média. Calibração: $5.5\times 60k = 0.05$ dá $k =
\qty{1.5e-4}{}$ por (mmol/L)(dia); a \qty{10}{mmol/L}: \qty{9.1}{\%}. Com
hemácias vivendo 60 dias a idade média é de 30 dias e a HbA1c fica pela
metade para a mesma glicose: um paciente a \qty{10}{mmol/L} leria
\qty{4.5}{\%}, normal.
\end{solution}

\begin{solution}{pb:b3:endocrinology:1}
\textbf{1.} \qty{90}{mg/dL} $= \qty{0.9}{g/L} = 0.9/180 = \qty{5.0}{mmol/L}$;
$0.9\times 15 = \qty{13.5}{g}$ no volume de distribuição.
\textbf{2.} $75/15 = \qty{5}{g/L}$: uma elevação de \qty{500}{mg/dL},
\qty{28}{mmol/L}. O pico real é bem mais baixo porque a absorção se
espalha por uma hora enquanto a eliminação retira glicose à medida que
ela chega, e o fígado capta boa parte já na primeira passagem.
\textbf{3.} $\qty{75}{g}/\qty{60}{min} = \qty{1.25}{g/min}$; em
\qty{15}{L}: \qty{8.3}{mg/dL/min}, quatro vezes $u_{0}$.
\textbf{4.} Excesso estacionário $u/a = 8.3/0.02 = \qty{417}{mg/dL}$,
constante de tempo $1/a = \qty{50}{min}$; depois de uma hora, $417(1 -
\mathrm{e}^{-1.2}) = \qty{291}{mg/dL}$ acima do jejum: cerca de
\qty{380}{mg/dL}, \qty{21}{mmol/L}.
\textbf{5.} $L = 0.2\times 0.05/(0.02\times 0.1) = 5$; excesso
estacionário $417/6 = \qty{69}{mg/dL}$, seis vezes menor que sem
\omterm{def:b3:endocrinology:glucose}{insulina}.
\textbf{6.} O encéfalo queima \qty{120}{g} de glicose por dia, não
armazena nenhuma e pouco mais consegue usar: uma glicose baixa dá
confusão e coma em minutos. Uma glicose alta glica proteínas e, ao longo
dos anos, danifica os pequenos vasos da \omterm{def:b3:sensory-systems:retina}{retina}, do rim e dos nervos e as
grandes artérias.
\textbf{7.} $\lambda^{2} + (a + \gamma)\lambda + (a\gamma + s\beta) = 0$:
traço $-0.12$, determinante $0.002 + 0.01 = 0.012$, discriminante
$0.0144 - 0.048 = -0.0336$.
\textbf{8.} $\omega = \sqrt{0.0336}/2 = \qty{0.092}{min^{-1}}$; período
\qty{68}{min}; tempo de decaimento $2/0.12 = \qty{16.7}{min}$.
\textbf{9.} $\dot g(0) = 100(-0.06 + c\,\omega) = -2$, de modo que
$c\,\omega =
0.04$ e $c = 0.44$.
\textbf{10.} $g = 0$ quando $\tan\omega t = -1/c = -2.29$, $\omega t =
\pi - 1.16 = 1.98$, $t = \qty{21.6}{min}$. No primeiro mínimo, perto de
\qty{32}{min}: $100\,\mathrm{e}^{-1.94}(\cos 2.97 + 0.44\sin 2.97) =
14.4\times(-0.91) \approx -\qty{13}{mg/dL}$, glicose de \qty{77}{mg/dL}.
\textbf{11.} Em $t = 0$, $\dot i = 100\beta > 0$: a \omterm{def:b3:endocrinology:glucose}{insulina} sobe
enquanto $g > 0$, e atinge o pico quando $\beta g = \gamma i$, o que
exige $g$ ainda positivo --- de modo que o pico vem antes de a glicose
chegar ao jejum, e depois que a glicose começou a cair (ela cai desde $t = 0$). No modelo o pico fica perto de \qty{12}{min}.
\textbf{12.} Com a entrada espalhada por uma hora o pico ocorre quando a
absorção e a eliminação se equilibram, mais baixo e mais tarde (30--60
min); o retorno espera o fim da absorção, daí as duas horas; o excesso
de \omterm{def:b3:endocrinology:glucose}{insulina} é menor para uma entrada gradual, de modo que a queda é
leve.
\textbf{13.} $L = 2.5$. Sadio, $g^{*} = (2/0.02)/6 = \qty{16.7}{mg/dL}$;
resistente, $100/3.5 = \qty{28.6}{mg/dL}$: uma elevação de cerca de
\qty{12}{mg/dL} (\qty{0.7}{mmol/L}), com glicose de jejum perto de
\qty{102}{mg/dL}, \qty{5.7}{mmol/L} --- o limiar da glicemia de jejum
alterada.
\textbf{14.} $i^{*} = \beta g^{*}/\gamma$: sadio, $8.3$; resistente,
\qty{14.3}{mU/L}; razão $1.7$: resistência à \omterm{def:b3:endocrinology:glucose}{insulina} compensada,
hiperinsulinemia com glicose quase normal.
\textbf{15.} $L = 0.1\times 0.01/0.002 = 0.5$; $g^{*} = 100/1.5 =
\qty{66.7}{mg/dL}$, \qty{50}{mg/dL} acima do estado sadio:
\qty{2.8}{mmol/L}, glicose de jejum de \qty{7.8}{mmol/L}
(\qty{140}{mg/dL}), acima do limiar diabético.
\textbf{16.} Discriminante $(a - \gamma)^{2} - 4s\beta = 0.0064 - 0.004 =
0.0024 > 0$: retorno monótono. $\lambda = -0.06 \pm 0.0245$: $-0.0355$ e
$-0.0845$ por minuto; constante de tempo mais lenta de \qty{28}{min},
contra \qty{16.7}{min} do decaimento sadio: a alça que falha retorna
mais devagar e nunca cai abaixo do jejum.
\textbf{17.} Glicose alta com \omterm{def:b3:endocrinology:glucose}{insulina} alta: resistência com compensação
parcial --- tipo~2; o valor de duas horas de \qty{13}{mmol/L} passa de
\qty{11.1}{}, de modo que é diabetes, e não mera tolerância diminuída.
\textbf{18.} Nenhum peptídeo~C significa nenhuma \omterm{def:b3:endocrinology:glucose}{insulina} endógena:
tipo~1. Sem \omterm{def:b3:endocrinology:glucose}{insulina} o músculo e a gordura não conseguem captar glicose,
a gordura é degradada a ácidos graxos e corpos cetônicos, a proteína
muscular é catabolizada para a gliconeogênese, e a glicose se perde na
urina com suas calorias: o paciente passa fome em meio à fartura.
\textbf{19.} $5\times 9/5.5 = \qty{8.2}{\%}$.
\textbf{20.} $1.5\,\mathrm{e}^{0.46\times 3} = 1.5\times 3.97 =
\qty{6.0}{mU/L}$. T$_{4}$ dentro da faixa, TSH acima dela:
hipotireoidismo subclínico, com a glândula falhando e a \omterm{def:b3:endocrinology:axes}{hipófise}
compensando.
\textbf{21.} $1.5\,\mathrm{e}^{0.46\times 8} = 1.5\times 39.6 =
\qty{59}{mU/L}$. Uma T$_{4}$ baixa com um TSH de apenas $0.3$ significa
que a \omterm{def:b3:endocrinology:axes}{hipófise} não está respondendo: a lesão é central (\omterm{def:b3:endocrinology:axes}{hipófise} ou
\omterm{def:b3:endocrinology:axes}{hipotálamo}), e não na tireoide.
\textbf{22.} Uma meia-vida: o nível cai a \qty{50}{\%}, e os sintomas se
instalam ao longo de semanas. O cortisol, com meia-vida de
\qty{80}{min}, desaparece em horas depois de uma dose esquecida, e um
estresse naquele dia não encontra defesa: a semana esquecida de tiroxina
é desconfortável; o dia esquecido de cortisol pode matar.
\textbf{23.} T$_{4}$ alta (o \omterm{def:b3:adaptive-immunity:antibody}{anticorpo} aciona a glândula à revelia da
retroalimentação); TSH suprimido pela T$_{4}$ alta; a relação log-linear
o leva a ficar indetectável --- em $T_{4} = 30$, $1.5\,\mathrm{e}^{-6.9}
= \qty{0.0015}{mU/L}$ ---, de modo que um TSH indetectável é o primeiro
sinal.
\textbf{24.} O nível de um \omterm{def:b3:endocrinology:hormone}{hormônio} informa o equilíbrio entre sua ordem
e sua glândula: TSH alto com T$_{4}$ baixa é uma glândula falha, e TSH
alto com T$_{4}$ alta é uma \omterm{def:b3:endocrinology:axes}{hipófise} desgovernada; \omterm{def:b3:endocrinology:glucose}{insulina} alta com
glicose alta é um corpo resistente, e \omterm{def:b3:endocrinology:glucose}{insulina} baixa com glicose alta é
uma ilhota destruída. Lida sozinha, uma T$_{4}$ de \qty{7}{pmol/L} ou
uma \omterm{def:b3:endocrinology:glucose}{insulina} de \qty{30}{mU/L} não diz onde está a falha; lida com seu
comandante, diz.
\textbf{25.} \omterm{thm:b3:endocrinology:loop}{Ganho de alça} $L = 5$; período de \qty{68}{min} e tempo de
decaimento de \qty{17}{min}; glicose de jejum da alça que falha,
\qty{7.8}{mmol/L}.
\end{solution}
